\chapter{Classifying schemes and diagnosing connected arguments}
\label{sec:p18-scheme-classification-diagnosis}

Calling a passage an argument from authority may leave its central
ambiguity unresolved. Is a manager directing an action, an engineer
offering an expert judgment, or a witness reporting what happened?
Calling another passage consequentialist may conceal a classification,
a prediction, a value judgment, and a recommendation that need separate
assessment. The useful classification identifies the inferential work
each step performs.

This chapter develops the complete classification and diagnostic
presentation of Macagno, Walton, and Reed's \emph{Argumentation Schemes:
History, Classifications, and Computational Applications},
\emph{IFCoLog Journal of Logics and Their Applications} 4(8), 2017,
pp.~2493--2556 \cite{MacagnoWaltonReed2017Schemes}.
Its proposed purpose-based tree, its bottom-up scheme clusters, and
its analysis through connected schemes are the operative method.
The article's historical classifications are retained as identified
comparisons. Its reports of software and experiments are dated reports,
not a demonstration that this chapter implements a classifier.

The classification is a proposed interpretive instrument. It can narrow
the plausible schemes, expose alternative readings, and guide further
questions. It does not guarantee one correct label for every text.
The authors explicitly leave a fuller integration with discourse
interpretation as further work (pp.~2547--2548).

\section{Keep purpose, material connection, and inference distinct}
\label{sec:p18-scheme-dimensions}

\begin{hbdefinition}{def:p18-scheme-three-dimensions}{Three dimensions of a reconstructed inference}
The \emph{pragmatic purpose} is what the speaker seeks to accomplish
with the inference: for example, recommend an action, predict an event,
or establish a classification.

The \emph{material} or \emph{semantic relation} connects the subject
matters of premise and conclusion: cause and effect, means and goal,
species and genus, part and whole, similarity, a source and its report,
or a consequence and its evaluation. Here \emph{material} does not mean
material implication.

The \emph{reasoning form} specifies how the premises are used to reach
the conclusion: deduction, defeasible rule application, induction,
analogy, or abduction. A rule application may itself use a particular
pattern, such as modus ponens. Identify its actual premise and
qualification before assigning that logical description.

A \emph{scheme} captures a characteristic combination of these
dimensions. A \emph{scheme instance} is a reconstructed argument that
supplies its required roles and contents. A \emph{scheme family}
groups related schemes; a \emph{cluster} can also contain schemes linked
by composition or opposition, not just by subtype.

In the diagnostic record, a \emph{subtype} adds conditions to a
more general scheme while retaining its relevant roles.
A \emph{component} is an inference that helps supply a premise
or intermediate conclusion in a larger argument.
A \emph{counterargument} offers a reason against a premise,
conclusion, or inferential connection. These three relations
must be recorded separately.
Premises are \emph{linked} when they jointly supply one
inferential step; independent reasons provide distinct routes
to the same conclusion.
\end{hbdefinition}

The source's Table~11, pp.~2518--2519, illustrates deduction with
definition/genus, cause-to-effect, consequences, and commitment;
induction with argument from example; and abduction with imperfect
signs, practical reasoning, and best explanation. These are illustrative
matchings, not exclusive assignments. Analogical, inductive, deductive,
and abductive reasoning can organize the same material relation in
different arguments (pp.~2495--2496, 2523).

\begin{hbexample}{def:p18-scheme-causal-forms}{One causal connection, several inferences}
In a fictional device, let \(C\) be the proposition that a relay
is energized and \(E\) the proposition that an indicator is lit.
Energizing the relay, with the circuit intact and power available,
lights the indicator.
Consider five different inferences.
\begin{enumerate}
\item The relay is energized, the circuit is intact, and power is
available; therefore the indicator lights. With the stated exceptionless
device rule, this is forward deduction. With only a typical-operation
rule, its conclusion must retain that defeasible qualification.
\item The indicator does not light; therefore the relay is not energized.
This is justified by contraposition only if the other required
conditions are established and the relevant conditional is strict.
A broken circuit supplies a counterexample to unqualified reasoning
from a merely typical causal connection.
\item The indicator lights; therefore the relay may be energized.
This is an explanatory hypothesis. A test mode that lights the indicator
without energizing the relay is a competing explanation. Observing \(E\)
does not deductively establish \(C\).
\item The relay is not energized; therefore the indicator does not light.
This needs an additional exclusion of other ways to light the indicator.
Test mode again defeats the inference. Calling it abductive does not
supply the missing exclusion.
\item A previous device had both an energized relay and a lit indicator;
this device has a lit indicator; therefore its relay may be energized.
This is a comparison or tentative generalization requiring relevant
similarity and support for the extrapolation. A single remembered
case does not establish a general causal rule.
\end{enumerate}
All five involve the same causal subject matter, while their inferential
directions and evidential obligations differ. This preserves the source's
point about cross-cutting dimensions without treating default rules
as automatically closed under contraposition.
\end{hbexample}

\section{The full purpose-and-means classification}
\label{sec:p18-scheme-purpose-tree}

The source's Figures~4--7, pp.~2520--2522, supply two first questions.
Does the conclusion assess the desirability of an action or establish
the acceptability of a judgment? Does its offered force come from the
subject matter itself or from a source's standing?

\begin{hbdefinition}{def:p18-scheme-internal-external}{Internal and external support}
\emph{Internal support} uses features or relations of the subject
matter: a definition, causal relation, expected outcome, or means to
a goal. \emph{External support} appeals to a source's authority,
knowledge, reliability, or a social practice. This is the article's
classification terminology; external support is not thereby irrelevant,
and internal support is not thereby reliable.

An \emph{action assessment} recommends or opposes a course of conduct.
A \emph{judgment} attributes occurrence, existence, a factual category,
or an evaluative property. Evaluating a proposed action can itself
contain intermediate judgments. Classify those steps separately rather
than forcing the whole passage into one exclusive branch.
\end{hbdefinition}

\begin{hbmethod}{def:p18-scheme-full-branch-map}{Branch questions and the families they retrieve}
\begin{description}
\item[External support for action: directive authority.]
Does the reason concern a source's power or legitimate role in directing
conduct? The source places practical authority here. Identify the role,
its domain, the actual direction, and why compliance is relevant.
An order does not by itself establish an empirical proposition.
\item[External support for action: popular practice.]
Is the reason that other people pursue that course?
The source's family is argument from popular practice.
Identify the relevant group, actual practice, and proposed reason for
conformity. Coordination may supply a relevant consideration; prevalence
alone does not establish that a practice is beneficial.
\item[External support for judgment: knowledge.]
Is acceptance based on a source's epistemic position?
Figure~5 lists expert opinion, position to know, popular opinion, and
\emph{ad hominem} arguments. Determine whether specialist competence,
direct access, or a group's opinion is doing the work.
Negative considerations about the source's competence can also
challenge this route.
\item[External support for judgment: reliability.]
Is the move directed at whether the source is trustworthy?
Figure~5 again lists \emph{ad hominem} arguments. Its repetition marks
two possible targets, knowledge and reliability, rather than a second
independent piece of evidence. A relevant challenge can weaken reliance
on a report; an unrelated personal insult does not do so, and weakening
a report does not establish its contrary.
\item[Internal support for action: means to a goal.]
Does the reason connect an action to a pursued end?
Figure~6 lists practical reasoning, values, and distress.
Identify the agent, goal, action, feasibility, and the claimed
means--end connection. A distress appeal can seek help by making
relief of a difficulty salient; establish whose goal or value makes
that difficulty a reason for the proposed action.
\item[Internal support for action: consequences.]
Does the reason concern what the action would bring about and how
those outcomes are evaluated? Figure~6 lists consequences, danger,
fear, threat, and sunk cost. Separate prediction from evaluation.
Identify whether harm is merely warned of, emotionally emphasized,
or conditionally promised by an agent. Determine whether past
expenditure bears on future outcomes or is simply irrecoverable.
\item[Internal judgment: prediction.]
Does the conclusion concern a future event or state?
Figure~7 places cause-to-effect arguments here. State the causal
premise, relevant circumstances, and uncertainty.
The source's broader occurrence/existence discussion also includes
present events and entities; this is not restricted to future tense.
\item[Internal judgment: retrodiction.]
Does the conclusion infer a past occurrence or an unobserved cause?
Figure~7 places sign and abductive arguments here.
State the observed indication and compare alternative explanations.
A conventional indication may require a convention rather than a
physical cause.
\item[Internal judgment: factual classification.]
Does the conclusion ascribe a factual property or category?
Figure~7 lists verbal classification, composition, division, and sign.
Supply the definition or classification criterion, the case facts,
or the relevant part--whole relation. A whole's property is not
automatically each part's property, or conversely.
\item[Internal judgment: value classification.]
Does the conclusion ascribe a positive or negative evaluative property?
Figure~7 lists classification and sign. Identify the value,
its application criterion, and whose hierarchy of values is invoked.
An observed act may be used as a sign of a disposition, which then
receives an evaluation; these are distinguishable steps.
\end{description}
\end{hbmethod}

These entries include every leaf family named in the source's main
tree. The tree retrieves candidates, while the particular scheme's
premises and critical questions determine a more specific assessment.
For several labels, such as distress and sunk cost, the article gives
a location but no complete independent scheme entry. Their original
scheme presentations are not acquired merely by citing the article.
The diagnostic examples below explain how the indicated distinctions
operate without attributing an invented universal premise list to
those names.

A conclusion's wording alone may be ambiguous. \emph{Unacceptable}
can be a value classification, a shorthand recommendation to prohibit,
or a report of an institution's decision. \emph{Should} can express
an expectation as well as a practical recommendation.
Use its context and role in the inference before taking a branch.

\section{A usable reconstruction and diagnosis procedure}
\label{sec:p18-scheme-diagnostic-procedure}

\begin{hbmethod}{def:p18-scheme-diagnosis-method}{Top-down selection with bottom-up reconstruction}
The inputs are a passage, its context, the issue, the participants'
attributable commitments, and available evidence. The output is a
record of inferential steps, candidate schemes, missing premises,
critical questions, and a qualified assessment.
\begin{enumerate}
\item State the intended conclusion with its actual scope and qualifier.
Separate a quoted speaker's claim from the current speaker's endorsement.
Record rival readings when the wording does not settle the purpose.
\item Identify the immediate premises offered for that conclusion.
Distinguish explicit premises, contextually attributable implicit
premises, and premises the argument would need but has not established.
Do not promote the third category into the second merely to make a
scheme fit.
\item Apply the action--judgment and internal--external questions.
For each surviving branch, write the answer and the textual reason
for it. Eliminate a candidate only when a required role is absent
or contradicted; a missing word alone need not eliminate an implicit role.
\item Refine within the complete branch map: directive role or practice;
knowledge or reliability; means or consequences; occurrence or
classification; factual or evaluative classification.
Retain parallel classifications when a step combines reasons.
\item Identify the material relation and the actual reasoning form
separately. For a causal relation, for example, distinguish forward
prediction, reverse explanation, and extrapolation from examples.
Preserve strict, typical, probabilistic, and hypothetical qualifications.
\item Instantiate the most specific supported scheme. Write the
agent, source, domain, action, goal, value, event, or category in its
proper role. Stop at a family-level diagnosis if the passage does not
support a narrower interpretation. Record what information would
distinguish the remaining candidates.
\item Inspect the candidate's cluster. Is another scheme a more specific
variant, a component supporting a premise, or a counterargument?
These relations answer different questions and should not share one
undifferentiated arrow.
\item Decompose the passage into a net when one scheme cannot account
for its steps. Assign an identifier to each proposition and each
inferential step; show which conclusion supplies which later premise.
Record linked premises, independent supporting reasons, and attacks
distinctly. Keep the provenance of every added premise visible.
\item Ask the critical questions actually associated with the selected
presentation. Specify their targets: premise, applicability or
exception, inference, conclusion, or competing argument.
For their burden and response effects, use the selected dialogue
procedure rather than inferring a universal default from the scheme name.
\item Give an assessment for each step and for the ultimate conclusion.
A successful classification is not confirmation of the premises,
an empirical causal finding, or a proof that the recommendation wins.
State unresolved alternatives, missing links, and relevant defeaters.
\item For argument construction, reverse the search from the intended
conclusion: choose a suitable branch, look for actual supporting
premises, and expose the resulting gaps. An available scheme suggests
what evidence to seek; it does not generate true evidence.
\end{enumerate}
This numbered workflow operationalizes the source's top-down and
bottom-up directions (\S\S2.3, 5--7). It is a manual interpretive
procedure, not a claimed total algorithm for unrestricted language.
\end{hbmethod}

\section{The decision-making cluster and its distinct schemes}
\label{sec:p18-scheme-practical-cluster}

\subsection{Instrumental reasoning and its five questions}

The article's Table~12, p.~2524, starts with an agent's goal \(G\)
and the claim that action \(A\) is a means to \(G\), then concludes
that the agent practically ought to perform \(A\).
This is a defeasible starting reason. It does not say that \(A\)
is the only means or the best means.
Its five critical questions, in independently worded form, are:

\begin{enumerate}
\item Which other goals could conflict with \(G\)?
\item Which alternative actions could also achieve \(G\)?
\item Which available means is most efficient for the relevant comparison?
\item What supports the practical feasibility of \(A\)?
\item Which further consequences of \(A\) should enter the assessment?
\end{enumerate}

The source also supplies five identification conditions
(pp.~2528--2529): an agent or group faces an action decision;
the circumstances supply pro and con considerations; goals and the
perceived circumstances guide the choice; those reasons are compared
for strength; and the agent presents that comparison as the basis
for acting or refraining. Several available actions may need comparison.
The value-based variant adds a sixth condition: the decision is
also justified by the agent's values.
These are proposed identification conditions, not observations that
every actual agent reasons this way.

\subsection{Consequences, values, and goal-based reasoning}

\begin{hbdefinition}{def:p18-scheme-consequence-value-forms}{Distinct members of the practical cluster}
The article's Tables~13--17, pp.~2525--2527, distinguish:
\begin{description}
\item[Positive consequences.] Performing \(A\) would plausibly bring
about favorably evaluated consequences; this provides a defeasible
reason to perform \(A\).
\item[Negative consequences.] Performing \(A\) would plausibly bring
about unfavorably evaluated consequences; this provides a defeasible
reason against \(A\).
\item[Positive value.] Agent \(s\) judges value \(V\) positive, and
its positive standing is a reason for \(s\) to commit to goal \(G\);
therefore \(V\) supplies that reason for commitment to \(G\).
\item[Negative value.] Agent \(s\) judges \(V\) negative, and that
negative standing is a reason to retract commitment to \(G\);
therefore \(V\) supplies a reason for that retraction.
\item[Goal-based practical reasoning.] The agent has goal \(G\);
its values support \(G\); and performing \(A\) is claimed to be
necessary or sufficient for \(G\); therefore there is a practical
recommendation to perform \(A\).
\end{description}
The bridge connecting a value to this particular goal is a premise,
not something obtained from the word \emph{positive} alone.
The last scheme's necessary and sufficient means variants differ:
necessity excludes achieving the retained goal without the means;
sufficiency permits the means as a route, without excluding alternatives.
Neither reading suppresses competing goals, feasibility, costs, or
the possibility of reconsidering \(G\).
\end{hbdefinition}

The article contrasts its separate value-to-goal reasoning with
presentations in which values are integrated into a transition from
current circumstances to a new situation. It refers to Atkinson,
Bench-Capon, and McBurney's 2005 formulation for the latter.
The referral identifies a distinct presentation; it is not a full
acquisition or formal derivation of that source here.

The cluster has several relationships (\S6, Figure~8).
Value-based practical reasoning adds value support to instrumental
reasoning. A negative-consequence argument can oppose the proposed
action or challenge the adequacy of its practical assessment.
Positive and negative consequence arguments depend on positive and
negative evaluation. The authors propose treating them as more specific
members of the corresponding value families.
These are their proposed classificatory relations, not an assertion
that every corresponding logical formula entails every other one.

\subsection{Slippery slopes: the sequence and the missing mechanism}

The article's basic slope scheme (Table~18, p.~2527) has an initially
plausible proposal \(A_0\), a claimed sequence in which it leads to
\(A_1\), then \(A_2\), and eventually \(A_n\), and an adverse evaluation
of \(A_n\). It offers a reason to reject \(A_0\).
Its three critical questions ask:

\begin{enumerate}
\item Which intervening propositions linking \(A_0\) to \(A_n\)
are actually supplied?
\item Which additional steps must be supplied to make the sequence
plausible?
\item Which are the weakest links, where the claimed transition
especially needs examination?
\end{enumerate}

The later discussion, pp.~2529--2530, emphasizes three distinguishing
elements: recursive transitions, a grey zone in which a stopping
boundary is uncertain, and a loss-of-control mechanism.
The \emph{grey zone} concerns an alleged inability to draw or maintain
a relevant stopping distinction; \emph{loss of control} concerns
why the agent would not stop the progression before its bad outcome.
Those premises need reconstruction and support. A long list of events
alone is not evidence for either.

The article identifies causal, precedent, and linguistic-vagueness
variants, plus an all-in combination, and proposes a minimal basic
scheme underlying specialized forms. In a causal variant, the links
claim event production. In a precedent variant, one decision is
claimed to commit or pressure later decisions through relevant similarity.
In a linguistic variant, an unstable boundary in a term or criterion
is claimed to prevent principled stopping. An all-in argument combines
such mechanisms. The article's reference to Walton's 1992 treatment
does not supply all of that book's original variants.

A slope belongs to the negative-consequence family and also draws
on practical goals and values. Figure~8 links the precedent slope
to both precedent reasoning and the basic slope.
This is a cluster with multiple relations, not a disjoint partition
in which a passage must receive exactly one label.

\section{Worked diagnoses across the branches}
\label{sec:p18-scheme-diagnosis-cases}

All cases in this section are constructed for diagnosis.
Their stipulated facts are not empirical observations.
Each case identifies the route, the reconstruction, its unresolved
questions, and the resulting assessment.

\begin{hbexample}{def:p18-scheme-authority-diagnosis}{One source, three different roles}
An archive committee hears three statements:
\begin{enumerate}
\item \emph{The director has instructed us to close Room B tonight,
so we should close it.}
\item \emph{The director is a qualified ventilation engineer and
reports that Room B's fan is defective, so its fan is probably defective.}
\item \emph{The director watched the fan stop at noon and reports
that observation, so it probably stopped at noon.}
\end{enumerate}
The first conclusion is an action recommendation supported externally
by a directive role. Its needed premises concern the instruction's
authenticity, the director's authority over this action, and relevant
limits on that authority. It does not establish a fan defect.

The second concludes a judgment and uses specialist competence:
external support through knowledge, specifically expert opinion.
Instantiate source, technical field, reported proposition, and
field relevance. Ask about expertise, reliability, agreement with
other competent sources, and the basis of the diagnosis.

The third also uses external knowledge, but its immediate warrant
is opportunity to observe: position to know or witness testimony.
Ask whether the director was present, could perceive the fan,
and reported accurately. Engineering competence may help interpret
the observation, but seeing a stopped fan does not itself establish
its cause. If the passage then infers a defect, add a separate
abductive step. The same person therefore occupies different roles;
the word \emph{director} cannot select the scheme by itself.
\end{hbexample}

\begin{hbexample}{def:p18-scheme-popularity-diagnosis}{Popular practice, popular opinion, and reliability}
Suppose a committee argues: \emph{All partner archives use format F,
so we should use F.} Its conclusion is practical and its cited reason
is a group's practice. Reconstruct the additional proposed warrant:
using a shared format reduces conversion failures with those partners.
If the stipulated interfaces actually accept F and alternatives are
costlier, that coordination consideration supplies internal practical
support alongside the external practice. If the partners' systems
are unrelated to this archive's exchange needs, prevalence alone
does not establish that F is a suitable choice.

Now change the conclusion: \emph{All partner archivists think F is
lossless, so F is lossless.} The route is external knowledge through
popular opinion. The group claim is a report about opinion, and
the missing bridge is its evidential relevance to losslessness.
Popularity cannot replace the technical analysis. If qualified
specialists have independently tested the same property, expert
agreement and test evidence may form additional steps.

Finally, a critic says that the person reporting those tests sells F.
This is a potentially relevant reliability challenge if the financial
connection affects reporting or selective disclosure; it is not proof
that F loses data. A criticism of the person's hairstyle supplies
no such connection. The diagnostic output distinguishes knowledge,
reliability, and an irrelevant personal attack instead of marking
every reference to a person as the same fallacy.
\end{hbexample}

\begin{hbexample}{def:p18-scheme-occurrence-diagnosis}{Predicting failure and explaining an alarm}
A technician predicts that a stopped coolant pump will raise a device's
temperature. The conclusion is a future-state judgment, supported
internally by a causal relation. Its reconstruction needs the
operating load, the pump's role, and whether backup cooling is active.
Without those conditions, the prediction remains qualified.

A different speaker observes a temperature alarm and infers that
the pump probably stopped. This is internal retrodiction or diagnosis
from sign through abduction. Failed calibration, a blocked channel,
and an unusually high load are alternative explanations.
The same causal vocabulary does not make the reverse inference
deductive. A record of pump speed could instead support the pump
claim through an observation/sign step and make the net more specific.

If the alarm is merely a maintenance reminder programmed for noon,
its evidential connection is conventional scheduling rather than
temperature causation. The word \emph{alarm} does not settle that issue.
Classify at the sign-family level until the indicator's operation is known.
\end{hbexample}

\begin{hbexample}{def:p18-scheme-classification-diagnosis}{Factual category, evaluation, and part--whole errors}
An exhibition's stated classification rule calls an object a
\emph{loan item} exactly when another institution retains ownership
and has temporarily transferred custody. The record establishes
both conditions for vase \(v\). Classifying \(v\) as a loan item
uses an internal definitional connection, not an expert's authority
merely because a curator utters the conclusion.

Calling the loan \emph{unfair}, however, requires an evaluative
criterion. Suppose the critic specifies that a loan demanding all
conservation costs from an institution unable to influence display
conditions distributes responsibilities unfairly, and supplies those
features. That is value classification. The criterion and its
application remain open to challenge; the factual loan definition
does not itself imply unfairness.

The statement \emph{Every crate is light enough for one person;
therefore the whole shipment is light enough for one person}
attempts composition. Ten stipulated crates of \(8\) kilograms
form an \(80\)-kilogram shipment. If the task's stipulated single-person
limit is \(15\) kilograms, the premises do not support the conclusion.
Conversely, an \(80\)-kilogram shipment need not make every crate
too heavy: that would be division. The transfer requires a
property-specific part--whole warrant.

If repeated disregard of documented handling restrictions is offered
as a sign of a careless disposition, identify the inference from acts
to disposition separately from the negative evaluation of carelessness.
A single accidental slip does not supply the same basis.
\end{hbexample}

\begin{hbexample}{def:p18-scheme-practical-diagnosis}{Goal, means, values, and alternatives}
A team says: \emph{We want visitors to find the temporary entrance.
Installing illuminated signs will achieve that, so we should install
them.} The route is internal action support through means to a goal.
Instantiate agent, goal, means, and feasibility; then apply all five
questions.

The team also seeks quiet evenings, so noise from installation could
conflict with another goal. Painted arrows are an alternative means.
If illuminated signs cost \(900\) credits and arrows \(150\), cost is
one efficiency consideration, but their visibility may differ.
Permission and electrical access bear on feasibility.
Glare affecting nearby rooms is a further consequence.
The initial practical reason therefore does not decide the comparison.

Now the team adds: \emph{Equal access matters, and visitors with limited
vision need high-contrast guidance.} This makes value support for
the goal explicit and prompts a classification of the alternatives'
accessibility. It yields a value-based practical net.
Positive evaluation of equal access alone would not justify an
unrelated goal: the value-to-goal bridge must be shown.

If illuminated signs are the only feasible means under the stipulated
constraints, their necessity is stronger than the initial claim that
they are one means. It still leaves the possible revision of the
goal, the constraints, and the cost assessment visible.
The article's Table~12 means formulation and Table~17
necessary-or-sufficient variants must not be silently collapsed.
\end{hbexample}

\begin{hbexample}{def:p18-scheme-distress-threat-diagnosis}{Distress, danger, fear, and threat}
Consider four requests concerning the same unlocked equipment store.
\begin{description}
\item[Distress.] A volunteer cannot carry vulnerable equipment to safety
alone and asks for help. The practical route is relief of a difficulty
as a goal. Needed premises concern the difficulty, the proposed help's
effectiveness, feasibility, and why aiding this person is a goal or
value of the addressee.
\item[Danger.] A supervisor warns that leaving the store unlocked
exposes equipment to theft. The route uses a predicted bad consequence.
Ask about actual access, the likelihood of loss, and the effectiveness
and costs of locking it.
\item[Fear.] The same warning vividly emphasizes the prospect of loss.
Emotional salience does not independently establish its probability.
The diagnosis still needs the consequence and evaluation premises.
\item[Threat.] A speaker announces that they will withdraw their
equipment loan unless the store is locked. The claimed consequence
is now an agent's conditional action. Its credibility depends on
that agent's control and willingness, and its practical relevance
on the loan's value. Those premises differ from an impersonal risk
of theft. The threatened loss does not establish that an unrelated
factual claim by the speaker is true.
\end{description}
The source places these families in practical branches. Their
classification identifies the offered reason; it does not endorse
the conduct or supply a universal moral assessment.
\end{hbexample}

\begin{hbexample}{def:p18-scheme-sunk-cost-diagnosis}{Past expenditure and future choice}
A team argues: \emph{We already spent \(100\) credits on the display,
so we must finish it.} The conclusion is practical and the cited
premise is sunk expenditure. Ask what future consequence links
that expenditure to continuing.

Suppose it is irrecoverable in either option. Continuing costs
another \(10\) and yields benefit \(8\); stopping yields neither.
On the stipulated monetary comparison, totals are \(-102\) for
continuing and \(-100\) for stopping. The common past \(100\)
does not make continuing preferable. If stopping instead triggers
an additional contractual cost, loses reusable work, or frustrates
a separately defended commitment, that is a new future-consequence
or commitment premise to assess. Do not confuse it with the
mere fact that money has already been spent.
\end{hbexample}

\begin{hbexample}{def:p18-scheme-slope-diagnosis}{Three slope mechanisms and a compressed warning}
A fictional archive considers a temporary exception permitting
an unlogged loan of one low-value object.
The critic gives the compressed warning that accepting it will
ultimately destroy the collection's audit trail.

The conclusion opposes an action, and the route is internal negative
consequences. To diagnose a slope, reconstruct the actual links:
\(A_0\), this exception; \(A_1\), subsequent staff treating it as
permission for comparable cases; \(A_2\), extension to cases whose
relevant limits are unclear; \(A_3\), inability to reconstruct
custody records. The badness of \(A_3\) rests on the archive's
accountability goal and valuation of traceability.

The precedent version requires reasons that this decision creates
a commitment or pressure to treat relevant later cases alike.
An explicit, enforceable exceptional criterion could block that link.
The linguistic version alleges a grey zone in \emph{low value}:
slightly more valuable objects are repeatedly treated as equivalent.
It needs evidence that a fixed threshold or review rule cannot
maintain a stopping boundary. Vagueness alone does not establish that.
The causal version claims that repeated unlogged movement itself
destroys reconstructible records; backup logs could block this link.
An all-in version combines those mechanisms, each of which retains
its own support requirement.

Apply the three slope questions: the critic supplies \(A_0,A_3\);
the reconstruction exposes the missing \(A_1,A_2\);
the weakest links concern extension of permission and inability
to maintain a boundary. Without support for those links and loss
of control, the short warning supports at most an incompletely
specified negative-consequence concern. A diagnosis must not
invent the intervening history and call it established evidence.
\end{hbexample}

\section{Reconstruct a net instead of forcing one label}
\label{sec:p18-scheme-nets}

The article's \S7, pp.~2531--2534, analyzes two compressed political
statements. Its first net links a classification of conduct to
attributed commitments and then adverse consequences. Its second
links a classification, value judgments, and an alleged escalation
sequence. Those are the authors' interpretations of historical
utterances, not fresh verification here of the events, attributed
commitments, or causal predictions.

The methodological lesson is to expose the intermediate propositions.
It does not authorize filling missing causal links with whatever
would make the favored interpretation work. Mark an interpretive
hypothesis as such, and distinguish it from an established premise.

\begin{hbexample}{def:p18-scheme-net-diagnosis}{A complete new net: a public image release}
A museum team argues: \emph{This draft includes donor addresses.
We promised to publish only anonymized material. Release would expose
donors and damage trust, so hold the draft and remove the addresses.}
The case uses a fictional internal promise and no assumed legal rule.

The main purpose is action guidance. Write the following propositions,
preserving their status.
\begin{description}
\item[\(P_1\), explicit.] The draft contains donor addresses.
\item[\(P_2\), needed classification premise.] A file containing such
identifying addresses is not anonymized under the team's stated
publication criterion.
\item[\(P_3\), intermediate.] This draft is not anonymized.
\item[\(P_4\), explicit.] The team committed to publish only anonymized
material.
\item[\(P_5\), intermediate.] The present draft fails that commitment's
release condition.
\item[\(P_6\), explicit prediction.] Releasing the addresses would expose
donors and damage their trust.
\item[\(P_7\), needed value premise.] The team regards the avoidable
exposure and loss of trust as adverse consequences worth preventing.
\item[\(P_8\), explicit practical conclusion.] Hold the present release.
\item[\(P_9\), needed means premise.] Removing the addresses alone
would satisfy the chosen anonymization criterion.
\item[\(P_{10}\), contextual goal.] Publish useful material consistently
with the promise.
\item[\(P_{11}\), explicit recommendation.] Revise by removing addresses
before publishing.
\end{description}

The inference record has five distinguishable components:
\begin{enumerate}
\item \(P_1,P_2\) support \(P_3\) by internal factual classification.
Verify the file content and obtain the actual criterion.
\item \(P_3,P_4\) support \(P_5\) by applying an attributable
commitment to the classified case. Its force depends on the promise's
actual terms and continued applicability.
\item \(P_6,P_7\) support \(P_8\) by negative consequences.
The prediction and adverse evaluation need different support.
\item \(P_5\) supplies an additional practical reason for \(P_8\),
given the team's commitment to act consistently with its promise.
Record that practical bridge rather than assuming every statement
of a promise already entails every action recommendation.
\item \(P_9,P_{10}\), together with the present defect, support
\(P_{11}\) by means--end reasoning. Check whether removing addresses
also removes other identifying information, whether revision is
feasible, and whether another way meets the goal better.
\end{enumerate}

The ultimate assessment is conditional. The text supplies the
promise and alleged addresses, while the reconstruction exposes
the anonymization criterion, value premise, and sufficiency of the
proposed revision. A demonstration that removing addresses alone
leaves identifiable faces attacks the sufficiency of \(P_9\).
It need not refute the negative-consequence reason to hold the draft.
Thus one failed revision proposal does not automatically validate
immediate release.
\end{hbexample}

\subsection{Comparison with the earlier scheme entries}

Chapter~\ref{sec:p18-argument-schemes-critical-questions} supplies
source-specific expert, witness, analogy, causal, and practical
entries with their particular questions.
The present method retrieves and combines those entries as follows.
\begin{description}
\item[Expert opinion.] External support for a judgment through knowledge.
Keep the specialist field and claimed proposition; do not replace
the earlier six-question source presentation with the tree's family name.
\item[Witness testimony.] Usually external support through position
to know, with observational and reporting conditions.
An explanatory inference from what the witness reports is an
additional internal step.
\item[Analogy.] A comparison-based reasoning form that can support
a factual classification, value judgment, causal hypothesis,
or practical conclusion. It crosses this purpose tree rather than
appearing as a single exclusive leaf. Preserve the earlier analogy
variants and their different similarity requirements.
\item[Causal inference and sign.] A material connection can organize
prediction, retrodiction, diagnosis, or practical consequences.
The inferential direction and alternatives determine the needed
assessment. A sign can also concern classification rather than cause.
\item[Practical reasoning.] The earlier necessary-condition presentation,
the present article's simple means formulation, and its value-based
necessary-or-sufficient formulation are distinct variants.
Retain the actual goal, means relation, and question list.
\end{description}

The commitment and challenge effects of the surrounding exchange
belong to Chapter~\ref{sec:p18-dialogue-burden}.
A scheme net does not by itself decide whether silence concedes
a premise or who must produce evidence in a dialogue.

\section{Historical classification dimensions retained by the article}
\label{sec:p18-scheme-classification-history}

The article's \S\S3--4, pp.~2500--2518, compares historical and modern
systems to motivate its own separation of dimensions.
The following is a map of that article's presentation.
It does not claim independent acquisition of Aristotle, Cicero,
Boethius, Abelard, or every modern work it cites.
The historical family names identify their roles in the comparison;
they are not interchangeable labels for the proposed tree above.

\subsection{Topics, maxims, and subject-matter relations}

The article reconstructs an Aristotelian \emph{topic} as a general
inferential principle that can warrant an enthymeme or help find
a more specific warrant. It distinguishes a generic topic from
domain-specific accepted premises. A general classification relation,
for example, differs from a particular disciplinary rule.
Finding a warrant through that relation does not establish that
the specific premise is true or accepted by the relevant audience.

Its Cicero map (Tables~3--4, p.~2503) separates intrinsic from
extrinsic support. Intrinsic support directly from the subject matter
includes definition by material parts or essential parts, and the
explanation of a name. Relations to other things include inflection,
genus, species, similarity, difference, opposites, concomitance,
antecedent, consequent, incompatibility, cause, effect, and comparison
of greater, lesser, or equal. Extrinsic support uses authority.
The separate issue classification asks whether the dispute concerns
conjecture, definition, or qualification. Cause, effect, and circumstances
help conjecture; descriptive and classificatory relations help
definition; comparison helps qualification. Thus issue type and
subject-matter relation are already distinguishable organizing concerns.

The Boethian account distinguishes a \emph{maxim}, a general
warranting principle, from its \emph{differentia}, the family of
relations under which such maxims are found.
For instance, a genus--species relation can provide an appropriate
classification warrant. If a property is predicated of every member
of a genus, it applies to members of its species; a property peculiar
to one species need not apply to its whole genus.
This direction must remain explicit.

The article's complete dialectical family map in Table~6, p.~2505,
has three branches:
\begin{description}
\item[Intrinsic.] From substance: definition, description, and explanation
of a name. From accompanying features: genus and species, integral
whole and its parts, efficient cause, matter, end, form, generation
or effect, corruption, uses, and associated accidents.
\item[Intermediate.] Inflections, coordinates, and division.
These concern grammatical or related semantic connections.
\item[Extrinsic.] Estimation of a thing, similarity, greater and lesser,
proportion, contraries, privation and possession, relative opposites,
affirmation and negation, and transumption.
The last is retained as the source's historical heading; the article
does not give an independent complete rule for every maxim beneath it.
\end{description}

Its rhetorical map (Table~7, p.~2506) concerns circumstances of
particular people and actions rather than exclusively conceptual
relations. The person branch includes name, nature, way of life,
fortune, studies, luck, feelings, disposition, purpose, deeds, and words.
The action branch includes the deed's nature, preceding and subsequent
events, what occurs during it, time and opportunity, place, manner,
and means. Comparison of circumstances uses species, genus,
contrary, result, greater, lesser, and equal.
Extrinsic considerations concern naming the deed, identifying its
actors or approving parties, relevant law/custom/agreement/judgment/
opinion/theory, conflict with custom, and general agreement.
These historical categories do not establish that a stereotype about
a person's origin or disposition is reliable evidence.

The Abelard discussion (pp.~2507--2508) separates the general
inferential maxim from the material assumption connecting the terms.
To classify an individual through a species--genus link, the
reconstruction needs both the individual's membership and the actual
inclusion of the species in the genus. Calling the relation a topic
does not supply that substantive premise. This distinction anticipates
the chapter's separate records for material connection and inference.

\subsection{The modern comparison map}

The modern systems show why classifications based on different purposes
can overlap. Preserve each system's organizing question.

\begin{description}
\item[Perelman and Olbrechts-Tyteca, as presented on p.~2509.]
Association and dissociation are the broad contrast.
Association contains three groups.
Quasi-logical arguments include contradiction/incompatibility,
identity/definition, analyticity/analysis/tautology, justice,
reciprocity, transitivity, part inclusion, division, comparison,
and sacrifice.
Arguments establishing a structure of reality use particular cases
(example, illustration, model and antimodel) or analogy/metaphor.
Arguments based on a structure of reality use sequential relations
(causal link, pragmatic consequence, ends and means, waste,
direction, unlimited development), coexistence (analogy, person
and acts, authority, speech as the speaker's act, group and members,
act and essence, symbolic relation), double hierarchy, and
differences of degree and order.
The article identifies mixed conceptual, logical, and relational
criteria; it does not claim these are one uniform dichotomy.

\item[Toulmin, Rieke, and Janik, p.~2510.]
The nine warrant-function groups in Figure~1 are analogy,
generalization, sign, cause, authority, dilemma, classification,
opposites, and degree. Some name a reasoning form, some a
subject-matter connection, and some a logical pattern.
The comparison motivates distinguishing those dimensions instead
of treating every warrant name as the same kind of category.

\item[Kienpointner, pp.~2510--2511.]
Four criteria are inference type, epistemic status of premises
(real or fictive), dialectical role (for or against a thesis),
and pragmatic role (descriptive or normative).
The top-level groups use a rule, establish a rule inductively,
or use or establish a rule.
Rule-using classification contains definition, genus--species,
and whole--part; comparison contains similarity, resemblance,
difference, and arguments from lesser or greater;
opposition contains contradictory, contrary, relative, and
incompatible relations; causal schemes contain cause, effect,
reasons, consequences, means, and ends.
Rule-establishing reasoning uses inductive examples.
The third group contains illustrative example, analogy, and authority.
These cross-cutting real/fictive, pro/con, and descriptive/normative
variants are not erased by choosing a content family.

\item[Pragma-dialectics, pp.~2511--2512.]
The three broad groups are symptomatic, similarity-based, and
instrumental or causal argumentation. The article places inherent
qualities, characteristic parts, and authority under the symptomatic
heading; analogy under similarity; and consequences or means--end
reasoning under the causal heading. Its criticism concerns mixing
a material relation with a reasoning type.
This historical grouping does not force every sign argument in
the present method to be a physical causal inference.

\item[Grennan, pp.~2512--2514.]
Nine warrant types are listed: cause to effect, effect to cause,
sign, sample to population, parallel case, analogy, population to
sample, authority, and ends--means.
The distinct eight claim types are obligation; supererogatory
action; prudential action; evaluation by grading, rating, or comparison;
physical-world fact, including institutional fact; mental-world
attribution; constitutive rule; and regulative rule.
The distinction between supererogatory and prudential action concerns
benefit to others versus the agent's own interests; constitutive
rules concern definitions or necessary connections, while regulative
rules state obligations or prohibitions.
The obligation illustration branches between obligation premises
and other premise kinds. The former has sample-to-population,
population-to-group, and effect-to-cause versions.
The latter has evaluative, physical-world, mental-world, and
regulative-rule premises; physical-world premises further divide
into institutional and brute-fact cases, while causal versions
appear for evaluative and regulative-rule premises.
Warrant type and conclusion type therefore form separate
classification dimensions.

\item[Katzav and Reed, pp.~2514--2517.]
A \emph{relation of conveyance} specifies how the fact expressed
by a premise is represented as necessitating or making liable
the fact expressed by a conclusion. The article's explicitly
nonexhaustive top-level map separates internal and external relations.
Internal relations comprise specification, constitution,
analyticity, and identity.
Specification includes species/genus and determinable/determinate.
Abstract constitution includes normative facts, with positive
and negative normative branches, nonnormative abstract facts,
necessary conditions, causal laws and singular causal conditionals,
constitution facts, possibility, and impossibility.
Concrete constitution includes kind instances, property instances,
properties constituted by properties or by particulars, singular
causal facts, a part's relation to its whole, and a whole's relation
to one of its parts.
Analyticity includes sameness of meaning, stipulative definition,
and implication; identity includes qualitative and numerical identity.
External relations divide into noncausal and causal dependence.
The noncausal branch includes noncausal laws, conservation of
quantity or quality, symmetry including spatial symmetry,
nomological incompatibility such as location/type
incompatibility, and topological structure.
The causal branch includes efficient cause, causal laws,
singular forward and reverse causal relations, common cause,
and final cause. The authors expressly deny a trivial one-to-one
mapping between these relations and the later scheme catalogue.

\item[Lumer and Dove, pp.~2517--2518.]
The three broad groups are deductive, probabilistic, and practical.
Deductive forms include elementary deduction and analytical forms,
including definition and legal subsumption.
Probabilistic forms include pure statistical/sign reasoning and
impure forms such as best explanation.
Practical forms include pure evaluation, justification of actions
or instruments, evaluation by adequacy conditions, welfare-ethical
value judgments, and practical arguments for theoretical theses.
The article observes that the first two broad groups concern
reasoning type while the third concerns practical purpose.
Identifying that mixture does not validate an inference from
desirable consequences to factual truth.
\end{description}

This comparison map explains the proposed tree's design.
It is not an acquisition of every original scheme beneath those
headings, a proof of the authors' historical interpretation, or
a finding that competing classifications are empirically inferior.

\section{Computational representation and the evidence boundary}
\label{sec:p18-scheme-computational-boundary}

\subsection{A proposal's rationale differs from its validation}

In \S2.2, pp.~2497--2499, the article describes an inductive
research program for justifying scheme recognition:
outline a scheme from the literature; analyze a substantial set
of examples and revise the scheme where needed; show the form's
importance in the relevant discourse; and provide empirical
justification for recognizing it as a basic scheme.
The article also invokes a teleological rationale: a defeasible
scheme may help an agent pursue epistemic or practical goals
better than foregoing an inference or using available alternatives,
especially when collecting more evidence is costly.
That rationale concerns a comparative use under circumstances.
It is not a proof that every scheme instance is true, that a
particular classifier is accurate, or that a teaching intervention
has a measured effect. Those are separate evidential questions.

\subsection{Recognition, evaluation, and construction have different outputs}

The article's \S8, pp.~2534--2539, distinguishes scheme structure
from the operational treatment of its critical questions.
A scheme can identify premise types and a conclusion type, so that
a candidate instance is checked for the required roles.
That structural fit does not establish its premises or its
acceptability against counterarguments.

The discussion of Verheij distinguishes four critical-question roles:
questioning a premise, identifying an exception, specifying proper-use
conditions, and identifying an attacking argument.
A premise attack, conclusion attack, and undercutting attack target
different components. An attacking argument might itself need
support. Preserve the selected presentation's treatment of questions
as assumptions, exceptions, or supported challenges rather than
turning every question into an automatically false premise.

The existing structured-argumentation construction in
Chapter~\ref{sec:p18-structured-argumentation}
supplies the exact strict/defeasible and attack/defeat distinctions.
Its abstract semantics must also be preserved: absence of an
accepted attacker is not by itself a complete rule for labeling an
argument accepted. In complete labeling, acceptance requires all
attackers to be out; attack by an in argument gives out; the remaining
cases can be undecided. A source's brief graph description does
not replace that full semantic interface.

The article describes formalized legal precedent schemes that
compare shared pro and con factors in a current case and a precedent,
using an appropriate preference among those factors. It also lists
ordinary meaning, technical meaning, precedent, purpose,
\emph{a contrario}, historical interpretation, and nonredundancy as
families of statutory interpretation arguments.
Their presence in this survey does not supply current legal rules,
all their defeat conditions, or the original formalizations.

For Carneades, the article distinguishes an earlier account using
ordinary premises, assumptions, and exceptions from the CAS2 model
implemented in version~4 at the time of writing.
The latter gives schemes roles in weighing, validating, and
generating arguments, and separates statements, arguments, and
issues in its graphs. An issue combines alternative positions
with its proof standard. The article contrasts version~3's backward
search with version~4's forward construction.
These are dated descriptions of identified systems.
No executable implementation, test of those systems, or complete
CAS2 semantics is supplied by naming their functions here.

\subsection{What the 2017 article actually reports about data and accuracy}

The reports in \S9, pp.~2539--2544, concern particular historical
datasets and studies. Keep their tasks and units distinct.

\begin{description}
\item[Internet Argument Corpus.]
The article reports \(390{,}000\) examples primarily organized
as quote--response pairs with polarity, with additional annotations
on subsets. It criticizes the gap between that unit and richer
argument structures. The number is not a count of fully analyzed
instances of every scheme.
\item[Potsdam Microtext Corpus.]
The article reports \(130\) human-written examples with constrained
five-component structures and English/German versions. It identifies
the fixed structure and small size as limitations on transfer to
unrestricted arguments. Parallel languages do not create independent
new examples for assessing generalization.
\item[Scheme-focused collections.]
It reports a Moral Maze collection with \(35\) scheme instances
and an ExpertOpinion--PositiveConsequences collection with \(71\)
examples. These restricted collections are not representative samples
of every argument family in ordinary discourse.
\item[Feng and Hirst classification.]
The article reports use of the 2008 catalogue's \(65\) schemes
and emphasis on five frequent families in the selected corpus:
example, cause to effect, practical reasoning, consequences,
and verbal classification. Those five accounted for a reported
\(61\%\) of the identified arguments in that database.
It reports classification accuracies ranging from \(0.64\) to \(0.98\).
This range is not a single accuracy for unrestricted end-to-end
argument reconstruction.
\item[Lawrence and Reed structure detection.]
The reported ranges are F1 \(0.59\)--\(0.91\) for components
and \(0.62\)--\(0.88\) for scheme instances.
These concern different outputs from classification given an
already identified argument. F1 and accuracy are also different
metrics; their numbers cannot be pooled as one measure.
\end{description}

These figures are reported by the 2017 article. The underlying
studies, datasets, annotation protocols, splits, and calculations
have not been independently reproduced by this chapter.
No new effectiveness estimate follows from the worked cases.
The article's claims about the future promise or limits of statistical
methods are its research assessment at that time, not a timeless
impossibility result or a current benchmark.

Its reported methodological idea remains useful to distinguish.
Lexical and semantic cues can suggest typed roles: reported speech
may indicate a source premise; a nearby claim of expertise may
support an expert-opinion hypothesis; a similar proposition may
be its candidate conclusion. Such cues constrain a search.
The word \emph{said} does not establish expertise, endorsement,
truth, or even the existence of an argument.

\begin{hbexample}{def:p18-scheme-mining-diagnosis}{A keyword match that needs interpretation}
Consider: \emph{You say the report records that the expert said the
device was reliable, but the report contains no such statement.}
A keyword system may find both \emph{expert} and \emph{said}.
The passage, however, challenges the attribution rather than
endorsing reliability on the expert's authority.

The diagnostic output first distinguishes the reported attribution
from the speaker's actual objection. If an earlier speaker relied
on that attribution, this is a challenge to a source-report premise.
Without that context it is a disputed report, not a fully instantiated
expert-opinion argument whose conclusion is reliability.
Recognizing component types helps locate the ambiguity; it does
not resolve scope, negation, or the discourse role automatically.
\end{hbexample}

\subsection{Ontology structure, subtype, and inherited questions}

The article's \S10, pp.~2544--2546, describes AIF as an
argument-representation ontology with extension mechanisms.
Its expert-opinion snippet, Figure~11, distinguishes a knowledge-position
conclusion, field-expertise and knowledge-assertion premises,
credibility and backup-evidence presumptions, and inconsistency
or lack-of-reliability exceptions. It also represents the conflict
associated with an expert-inconsistency exception.
These are typed roles in the displayed ontology; the word
\emph{exception} does not assert that the exception actually holds.

The article identifies three uses of ontology reasoning.
First, it can trace dependencies through a graph. If an argument
uses \(X\) to support \(Y\), and another uses \(Y\) to support \(Z\),
then the latter conclusion \(Z\) depends along that path on \(X\).
This is a structural dependency, not proof that every composed
defeasible inference is an exceptionless entailment.
Second, an encoded subtype relation can classify a more specific
scheme, such as the described fear-appeal example within negative
consequences. Third, an appropriate subtype relation can make a
general scheme's questions available for a more specific instance.
All three depend on the actual ontology and its definitions.

This does not justify copying every question across every relation
in a cluster. Component-of, counterargument-to, and subtype-of
are different. For example, a negative-consequence objection to
practical reasoning is not for that reason an instance of every
practical-reasoning premise pattern.
Nor is the printed ontology snippet a complete standalone
specification of the full ontology and its reasoner.

The article also describes OVA/OVA+, AIFdb and the Argument Web
as analysis and corpus infrastructure, and Avicenna as an implemented
ontology-based identification application. It discusses argument
construction by searching commitments and schemes toward an
ultimate claim. Such reported implementations differ from the
authors' proposed classificatory rationale and from an observed
performance result. Their existence in the historical report
does not establish their present availability or install them here.

Its historical discussion also identifies ArguMed as a mapping
implementation and Araucaria as earlier analysis and corpus
infrastructure. Its discussion of the then-new Watson Debater
describes retrieval of pro and con claims while observing that
that tool did not then use schemes. Thus the article itself
does not equate every argument-mining or invention system with
an implementation of its classification tree.

\section{Using the diagnosis without overstating its result}

A complete diagnosis states what the passage is doing, shows
its explicit and reconstructed premises, identifies the material
relation and reasoning form, and records the applicable scheme
variants and challenges. It can end with several supported readings
or an unfilled premise. That is a usable outcome when the text
does not justify a unique label.

For the practical cluster, keep values, goals, means, consequences,
and the alleged slope mechanism distinct. For source-based judgments,
keep role authority, epistemic access, expertise, reliability,
and social prevalence distinct. For a net, evaluate the individual
steps as well as the ultimate recommendation.
The classification gives an organized search for reasons and
defeaters; its output is not itself evidence for those reasons.

\begin{proofboundary}{def:p18-scheme-classification-next-boundary}{Completed presentation and separate originals}
The supplied source-specific extent comprises the article's complete
purpose-and-means tree, decision-making cluster, scheme-net method,
historical comparison branches, and computational/empirical report
with its qualifications. The new cases perform the diagnosis
across all major branches and make their missing premises explicit.

The original 2008 scheme catalogue, Walton's 1992 slope variants,
the distinct Atkinson--Bench-Capon--McBurney transition presentation,
the original mining experiments, and the full AIF/CAS2/precedent
formalizations remain separate primary-source developments.
This article's references to them do not constitute their acquisition
or completion. Full empirical validation of a classifier would
also require an identified implementation, task, corpus,
annotation method, evaluation design, and reported results.
\end{proofboundary}
